% Generated by scripts/assemble.py; corrections belong in source/corrections.json.
\examhead{63}{重庆大学 2025 年数学分析试卷}{62}
% p063-b21; source PDF page 63
\par\Needspace{5\baselineskip}
\blockanchor{p063-b21}
\majorquestion{一}(20 分) 计算函数极限 $\lim_{x \to +\infty} x \left[ \frac{1}{e} - \left( \frac{x}{1 + x} \right)^x \right]$.\par
% p063-b22; source PDF page 63
\par\Needspace{5\baselineskip}
\par\Needspace{6\baselineskip}
\blockanchor{p063-b22}
\majorquestion{二}(20 分) 设 $S$ 为 $\mathbb{R}^3$ 中的光滑封闭曲面, 取外侧, 考虑第二类曲面积分\par
\nopagebreak[4]
% p063-b23; source PDF page 63
\blockanchor{p063-b23}
\displayformula{I = \iint_ {S} \left(x ^ {3} + y ^ {2025}\right) \mathrm{d} y \mathrm{d} z + \left(2 y ^ {3} + y + z ^ {2024}\right) \mathrm{d} z \mathrm{d} x + \left(3 z ^ {3} - 4 z - 1929 x\right) \mathrm{d} x \mathrm{d} y.}
% p064-b00; source PDF page 64
\blockanchor{p064-b00}
(1) 试确定曲面 $S$ 的方程, 使得积分 $I$ 的值最小, 并求出这个最小值.\par
% p064-b01; source PDF page 64
\blockanchor{p064-b01}
(2) 将第 (1) 中得到的曲面 $S$ 在第一卦限的部分记作 $S_{1}$, 求 $S_{1}$ 的切平面, 使得该切平面与三个坐标轴围成的几何体的体积最小.\par
% p064-b02; source PDF page 64
\par\Needspace{5\baselineskip}
\blockanchor{p064-b02}
\majorquestion{三}(20 分) 解答如下问题:\par
% p064-b03; source PDF page 64
\blockanchor{p064-b03}
(1) 已知数列 $\{x_{n}\}_{n=1}^{\infty}$ 满足 $|x_{n+1}-x_{n}|<r|x_{n}-x_{n-1}|$, $n=2,3,\cdots$. 其中 $0<r<1$ 为某个常数, 证明: 数列 $\{x_{n}\}_{n=1}^{\infty}$ 收敛.\par
% p064-b04; source PDF page 64
\blockanchor{p064-b04}
(2) 如果将第 (1) 小题中的常数 $r$ 换成 1, 是否仍有 $\{x_{n}\}_{n=1}^{\infty}$ 收敛? 证明你的结论.\par
% p064-b05; source PDF page 64
\par\Needspace{5\baselineskip}
\blockanchor{p064-b05}
\majorquestion{四}(10 分) 设函数 $f(x)$ 在 $[0, 1]$ 上具有连续的一阶导函数.\par
% p064-b06; source PDF page 64
\par\Needspace{6\baselineskip}
\blockanchor{p064-b06}
(1) 证明: 对任意的 $x \in [0,1]$, 有\par
\nopagebreak[4]
% p064-b07; source PDF page 64
\blockanchor{p064-b07}
\displayformula{| f (x) | \leq \int_ {0} ^ {1} | f (t) | \mathrm{d} t + \int_ {0} ^ {1} | f ^ {\prime} (t) | \mathrm{d} t.}
% p064-b08; source PDF page 64
\par\Needspace{6\baselineskip}
\blockanchor{p064-b08}
(2) 证明不等式:\par
\nopagebreak[4]
% p064-b09; source PDF page 64
\blockanchor{p064-b09}
\displayformula{\left| f \left(\frac {1}{2}\right) \right| \leq \int_ {0} ^ {1} | f (t) | \mathrm{d} t + \frac {1}{2} \int_ {0} ^ {1} | f ^ {\prime} (t) | \mathrm{d} t.}
% p064-b10; source PDF page 64
\par\Needspace{5\baselineskip}
\blockanchor{p064-b10}
\majorquestion{五}（15分）判断下列命题的正误，正确的需给出证明，错误的要举出反例。\par
% p064-b11; source PDF page 64
\blockanchor{p064-b11}
(1) 若级数 $\sum_{n=1}^{\infty} x_n$ 收敛, 则级数 $\sum_{n=1}^{\infty} (-1)^n \frac{x_n}{n}$ 也收敛.\par
% p064-b12; source PDF page 64
\blockanchor{p064-b12}
(2) 若连续函数列 $\{f_n(x)\}_{n=1}^{\infty}$ 在 $[0,1]$ 上逐点收敛于连续函数 $f(x)$, 则 $\{f_n(x)\}_{n=1}^{\infty}$ 在 $[0,1]$ 上一致收敛于 $f(x)$.\par
% p064-b13; source PDF page 64
\blockanchor{p064-b13}
(3) 若函数 $f(x), g(x)$ 都在 $[0, +\infty)$ 上一致连续，且 $g(x)$ 以 $x$ 轴为水平渐进线，则函数 $f(x)g(x)$ 在 $[0, +\infty)$ 上一致连续.\par
% p064-b14; source PDF page 64
\par\Needspace{5\baselineskip}
\blockanchor{p064-b14}
\majorquestion{六}(20 分) 设 $D$ 是 $\mathbb{R}^2$ 中由光滑简单封闭曲线 $C$ 所围成的闭区域, 二元函数 $f$ 和 $g$ 在 $D$ 上具有连续的二阶偏导数, 记 $\Delta f = \frac{\partial^2 f}{\partial x^2} + \frac{\partial^2 f}{\partial y^2}$.\par
% p064-b15; source PDF page 64
\par\Needspace{6\baselineskip}
\blockanchor{p064-b15}
(1) 证明等式:\par
\nopagebreak[4]
% p064-b16; source PDF page 64
\blockanchor{p064-b16}
\displayformula{\iint_ {D} (g \Delta f) \mathrm{d} x \mathrm{d} y = \int_ {C} g \frac {\partial f}{\partial \boldsymbol {n}} \mathrm{d} s - \iint_ {D} [ \operatorname{grad} (f) \cdot \operatorname{grad} (g) ] \mathrm{d} x \mathrm{d} y.}
% p064-b17; source PDF page 64
\par\Needspace{6\baselineskip}
\blockanchor{p064-b17}
其中 $\pmb{n}$ 为曲线 $C$ 的外法单位向量, $\frac{\partial f}{\partial n}$ 表示 $f$ 沿 $\pmb{n}$ 方向的方向导数, 向量\par
\nopagebreak[4]
% p064-b18; source PDF page 64
\blockanchor{p064-b18}
\displayformula{\operatorname{grad} (f) = \left(\frac {\partial f}{\partial x}, \frac {\partial f}{\partial g}\right), \operatorname{grad} (g) = \left(\frac {\partial g}{\partial x}, \frac {\partial g}{\partial g}\right).}
% p064-b19; source PDF page 64
\par\Needspace{6\baselineskip}
\blockanchor{p064-b19}
(2) 证明: 若 $f$ 在曲线 $C$ 上满足 $f \equiv 0$, 则\par
\nopagebreak[4]
% p064-b20; source PDF page 64
\blockanchor{p064-b20}
\displayformula{\iint_ {D} \left(| f | ^ {2} + | \Delta f | ^ {2}\right) \mathrm{d} x \mathrm{d} y \geq 2 \iint_ {D} | \operatorname{grad} (f) | ^ {2} \mathrm{d} x \mathrm{d} y.}
% p064-b21; source PDF page 64
\par\Needspace{5\baselineskip}
\blockanchor{p064-b21}
\majorquestion{七}(25 分) 解答如下问题:\par
% p064-b22; source PDF page 64
\blockanchor{p064-b22}
(1) 证明: 含参变量 $y$ 的广义积分 $\int_{0}^{+\infty} e^{-xy} \frac{\sin x}{x} \, \mathrm{d}x$ 在 $[0, +\infty)$ 上一致收敛.\par
% p064-b23; source PDF page 64
\blockanchor{p064-b23}
(2) 证明: 在 $(0, +\infty)$ 上, 含参变量 $y$ 的广义积分 $\int_{0}^{+\infty} e^{-xy} \sin x \, \mathrm{d}x$ 不一致收敛, 但内闭一致收敛.\par
% p064-b24; source PDF page 64
\blockanchor{p064-b24}
(3) 利用第 (1) 和第 (2) 小题中的结论计算 Dirichlet 积分 $\int_{0}^{+\infty} \frac{\sin x}{x} \, \mathrm{d}x$.\par
% p065-b00; source PDF page 65
\blockanchor{p065-b00}
(4) 利用第 (3) 小题中的结论证明: $\int_{0}^{+\infty}\left[\int_{0}^{+\infty}\frac{\sin x\sin(x+y)}{x(x+y)}\mathrm{d}y\right]\mathrm{d}x=\frac{\pi^{2}}{8}$.\par
% p065-b01; source PDF page 65
\par\Needspace{5\baselineskip}
\par\Needspace{6\baselineskip}
\blockanchor{p065-b01}
\majorquestion{八}(20 分) 设 $f(x)$ 为 $[- \pi, \pi]$ 上的实值 Riemann 可积函数. 给定区间 $I \subset [-\pi, \pi]$, $I$ 上的特征函数定义为\par
\nopagebreak[4]
% p065-b02; source PDF page 65
\blockanchor{p065-b02}
\displayformula{\chi_ {I} (x) = \left\{ \begin{array}{l l} 1, & x \in I; \\ 0, & x \notin I. \end{array} \right.}
% p065-b03; source PDF page 65
\par\Needspace{6\baselineskip}
\blockanchor{p065-b03}
如果  $I_{1}, I_{2}, \cdots, I_{N}$  为  $[-\pi, \pi]$  中两两不交的子空间， $a_{1}, a_{2}, \cdots, a_{N}$  均为常数，且  $\bigcup_{k=1}^{N} I_{k} = [-\pi, \pi]$ ，那么我们称有限线性组合\par
\nopagebreak[4]
% p065-b04; source PDF page 65
\blockanchor{p065-b04}
\displayformula{a _ {1} \chi_ {I _ {1}} + a _ {2} \chi_ {I _ {2}} + \dots + a _ {N} \chi_ {I _ {N}}}
% p065-b05; source PDF page 65
\blockanchor{p065-b05}
为 $[- \pi, \pi]$ 上的阶梯函数.\par
% p065-b06; source PDF page 65
\par\Needspace{6\baselineskip}
\blockanchor{p065-b06}
(1) 证明: 对任意的 $\varepsilon > 0$, 存在 $[- \pi, \pi]$ 上的阶梯函数 $\varphi(x)$ 和 $\psi(x)$ 使得\par
\nopagebreak[4]
% p065-b07; source PDF page 65
\blockanchor{p065-b07}
\displayformula{\varphi (x) \leq f (x) \leq \psi (x), x \in [ - \pi , \pi ].}
% p065-b08; source PDF page 65
\par\Needspace{6\baselineskip}
\blockanchor{p065-b08}
并且\par
\nopagebreak[4]
% p065-b09; source PDF page 65
\blockanchor{p065-b09}
\displayformula{\int_ {- \pi} ^ {\pi} | \varphi (x) - f (x) | \mathrm{d} x <   \varepsilon , \int_ {- \pi} ^ {\pi} | \psi (x) - f (x) | \mathrm{d} x <   \varepsilon .}
% p065-b10; source PDF page 65
\par\Needspace{6\baselineskip}
\blockanchor{p065-b10}
(2) 利用 (1) 中的结论证明\par
\nopagebreak[4]
% p065-b11; source PDF page 65
\blockanchor{p065-b11}
\displayformula{\lim _ {p \to + \infty} \int_ {- \pi} ^ {\pi} f (x) \sin (p x) \mathrm{d} x = \lim _ {p \to + \infty} \int_ {- \pi} ^ {\pi} f (x) \cos (p x) \mathrm{d} x = 0.}
% p065-b12; source PDF page 65
\par\Needspace{6\baselineskip}
\blockanchor{p065-b12}
(3) 将 $f(x)$ 延拓为 $(- \infty, + \infty)$ 上以 $2\pi$ 为周期的周期函数, 延拓之后的函数仍记为 $f(x)$. 证明: $f(x)$ 的 Fourier 级数的部分和 $S_{n}(x)$ 可写成\par
\nopagebreak[4]
% p065-b13; source PDF page 65
\blockanchor{p065-b13}
\displayformula{S _ {n} (x) = \frac {1}{\pi} \int_ {- \pi} ^ {\pi} \frac {f (x + t)}{2 \sin \frac {t}{2}} \sin \left(\frac {2 n + 1}{2} t\right) \mathrm{d} t, n = 1, 2, \dots .}
% p065-b14; source PDF page 65
\blockanchor{p065-b14}
(4) 在第 (3) 小题的条件下, 再假设 $x_0 \in (-\pi, \pi)$ 为 $f(x)$ 在 $[- \pi, \pi]$ 上的唯一间断点, 且 $f(x)$ 在 $x_0$ 处的两个单侧极限和两个单侧导数 (应该是导数的单侧极限) 都存在. 证明: $f(x)$ 的 Fourier 级数在 $x_0$ 点收敛于 $\frac{A + B}{2}$, 其中 $A, B$ 分别为 $f(x)$ 在 $x_0$ 点的左极限与有极限.\par
